\section{Case evidence and lexical indicator}\label{app:measures}

The \href{supplement.pdf}{supplementary case evidence} retains the full
rival register and five candidate-mechanism audits.
Whitespace counts describe text volume, not loaded context.
The evidence-boundary lexicon matches
\texttt{unmet}, \texttt{not accepted}, \texttt{not inferred},
\texttt{is inferred}, \texttt{remains open/unproved/unverified},
\texttt{not claimed}, \texttt{no ... is/was claimed/inferred}, and
\texttt{not established}, case-insensitively.
Density is matches per 1,000 whitespace-delimited entry words, not a
fraction of statements or a general hedging measure.
The day-one evidence-based ticking rule changed reporting conventions;
it is a rival explanation for the lexical increase. The matched sample
and qualified precision assessment are retained with revision metrics.
Appropriate caution is not evidence of failure or honesty.

Deferred intervention definitions and historical freeze hashes are in
the \href{supplement.pdf}{supplementary appendices}, not collected outcomes.
